\section{基本原则}
\begin{enumerate}
  \item 平等原则。考试就考一个问题，平等原则不要求在诉讼中当事人的权利义务相同。
        \begin{enumerate}
          \item 包括权利\zd{相同}。
          \item 也包括权利虽然不同，\zd{相对应}。
          \item \zd{绝对不能狭隘的把平等理解为相同}。
        \end{enumerate}
  \item 同等和对等原则。\zd{只适用于外国}当事人。
  \item 辩论原则。在\zd{诉讼}中，\zd{当事人}有权进行论证。针对的是\zd{事实和证据}。
        \begin{enumerate}
          \item 时间：\zd{贯穿民事诉讼的始终}。
          \item 适用范围：\zd{诉讼程序}。有纠纷你才能辩论，诉讼程序解决纠纷，需要辩论，\zd{其他程序不解决纠纷，不需要辩论}。
          \item 内容：实体和程序都可以。
          \item 形式：口头书面都可以。
        \end{enumerate}
  \item 处分原则。\zd{私权利主体}可以处分自已的私权利，针对的是\zd{诉讼请求层面}。
        \begin{enumerate}
          \item 在法定范围内。
          \item 考试常考的一种违反处分原则的做法，法院的\zd{判决超出原告的诉讼请求}。\qd{多管闲事}。
        \end{enumerate}
        \begin{kk}{约束性辩论原则}
          要让当事人的辩论，对法院的裁判形成约束。
          \begin{enumerate}
            \item 只有当事人提出并\zd{加以阻挡的事实}，法院才能予以认定。
            \item 对当事人都\zd{没有争议的事实}法院应当予以认定 \yjt \zd{受当事人自认约束} \yjt 辩论原则。
            \item 法院对\zd{证据}的调查原则上仅限于当事人提出的 \yjt 辩论原则、处分原则、平等原则。
          \end{enumerate}
        \end{kk}
  \item 诚信原则。\zd{约束所有}参与民事诉讼的主体。
  \item 检查监督原则。\zd{检察机关}作为我国的法律监督机关有权对\zd{民事诉讼活动}进行监督。
        \begin{enumerate}
          \item \zd{监督：把权力关进笼子里}。\zd{对象一定是公权力的行使}。法院 \yjt 审判、执行；
          \item 方式：
                \begin{enumerate}
                  \item 抗诉。只有一种抗诉\zd{再审抗诉}，不满足再审条件不能抗诉。
                  \item 检察建议。再审检察建议。\zd{一定要满足再审条件}；其他检察建议。
                \end{enumerate}
        \end{enumerate}
  \item 在线诉讼原则。谁同意对谁用，不同意各玩各的。
\end{enumerate}
